\documentclass[11pt,a4paper]{article}

% ----------------------------------------------------------
% Pacotes básicos
% ----------------------------------------------------------
\usepackage[utf8]{inputenc}
\usepackage[T1]{fontenc}
\usepackage[brazil]{babel}
\usepackage{lmodern}
\usepackage{geometry}
\usepackage{setspace}
\usepackage{indentfirst}
\usepackage{graphicx}
\usepackage{float}
\usepackage{booktabs}
\usepackage{array}
\usepackage{tabularx}
\usepackage{hyperref}
\usepackage{xcolor}
\usepackage{listings}
\usepackage{microtype}
\usepackage{titlesec}
\usepackage{caption}
\usepackage{enumitem}

% ----------------------------------------------------------
% Configurações de página - mesmo padrão do primeiro relatório
% ----------------------------------------------------------
\geometry{
    left=2.1cm,
    right=2.1cm,
    top=2.0cm,
    bottom=2.0cm
}

\setstretch{1.08}
\setlength{\parskip}{0pt}
\setlength{\textfloatsep}{8pt}
\setlength{\floatsep}{6pt}
\setlength{\intextsep}{6pt}
\titlespacing*{\section}{0pt}{0.9em}{0.4em}
\titlespacing*{\subsection}{0pt}{0.7em}{0.25em}
\captionsetup{font=small,skip=4pt}

% ----------------------------------------------------------
% Links
% ----------------------------------------------------------
\hypersetup{
    colorlinks=true,
    linkcolor=black,
    urlcolor=blue,
    citecolor=black,
    pdftitle={Equivalência Formal do Controlador de Vending Machine},
    pdfauthor={Luciano Antunes de Oliveira}
}
\urlstyle{same}

% ----------------------------------------------------------
% Blocos de código e relatórios
% ----------------------------------------------------------
\definecolor{codegray}{gray}{0.96}
\definecolor{codegreen}{RGB}{0,100,0}
\definecolor{codeblue}{RGB}{0,40,140}

\lstset{
    basicstyle=\ttfamily\footnotesize,
    breaklines=true,
    breakatwhitespace=false,
    frame=single,
    backgroundcolor=\color{codegray},
    columns=fullflexible,
    keepspaces=true,
    showstringspaces=false,
    numbers=left,
    numberstyle=\tiny\color{gray},
    keywordstyle=\color{codeblue}\bfseries,
    commentstyle=\color{codegreen},
    aboveskip=5pt,
    belowskip=5pt,
    literate=
      {á}{{\'a}}1 {à}{{\`a}}1 {ã}{{\~a}}1 {â}{{\^a}}1
      {é}{{\'e}}1 {ê}{{\^e}}1 {í}{{\'i}}1
      {ó}{{\'o}}1 {õ}{{\~o}}1 {ô}{{\^o}}1
      {ú}{{\'u}}1 {ç}{{\c{c}}}1
      {Á}{{\'A}}1 {É}{{\'E}}1 {Í}{{\'I}}1
      {Ó}{{\'O}}1 {Ú}{{\'U}}1 {Ç}{{\c{C}}}1
}

\lstdefinestyle{tclstyle}{
    language=tcl,
    morekeywords={set_svf,compile_ultra,read_db,read_sverilog,read_verilog,set_top,match,verify,report_status,report_svf_operation,report_matched_points,report_unmatched_points,report_passing_points,report_failing_points}
}

\lstdefinestyle{reportstyle}{
    language={},
    numbers=none,
    keywordstyle=\color{black},
    commentstyle=\color{black}
}

% ----------------------------------------------------------
% Imagens
% ----------------------------------------------------------
\graphicspath{{imagens/}}

\newcommand{\imagemrelatorio}[4]{
\begin{figure}[H]
\centering
\IfFileExists{imagens/#1}{
    \includegraphics[width=#2\textwidth]{imagens/#1}
}{
    \fbox{
    \begin{minipage}[c][4.5cm][c]{0.9\textwidth}
        \centering
        Imagem não encontrada: \texttt{imagens/#1}\\
        Insira a imagem correspondente na pasta \texttt{imagens/}.
    \end{minipage}
    }
}
\caption{#3}
\label{#4}
\end{figure}
}

\begin{document}
\pagenumbering{gobble}

% ----------------------------------------------------------
% CAPA
% ----------------------------------------------------------
\begin{titlepage}
    \centering

    {\large CI Expert -- Residência em Microeletrônica\par}
    {\large Trilha: RTL Design\par}
    {\large Módulo: Verificação de Equivalência Formal\par}

    \vspace{3cm}

    {\LARGE \textbf{Equivalência Formal do Controlador de Vending Machine}\par}

    \vspace{0.7cm}

    {\Large RTL versus Netlist com Guidance SVF\par}

    \vspace{1cm}

    {\large Relatório Técnico - Atividade 02\par}

    \vfill

    \begin{flushleft}
    \textbf{Aluno:} Luciano Antunes de Oliveira\\
    \textbf{Projeto:} Controlador digital de vending machine\\
    \textbf{Ferramentas:} Synopsys Design Compiler e Synopsys Formality\\
    \textbf{Repositório GitHub:} \url{https://github.com/luciano066/grupo_NN_vending}
    \end{flushleft}

    \vfill

    {\large 2026\par}
\end{titlepage}

\pagenumbering{roman}
\setcounter{page}{1}
\tableofcontents
\newpage
\pagenumbering{arabic}
\setcounter{page}{1}

% ----------------------------------------------------------
\section{Introdução}
% ----------------------------------------------------------

Este relatório apresenta a verificação de equivalência formal do controlador de vending machine desenvolvido em SystemVerilog na atividade anterior. O objetivo do fluxo é demonstrar que as netlists mapeadas produzidas pelo Synopsys Design Compiler preservam o comportamento sequencial do RTL original.

Diferentemente da simulação funcional, que avalia o circuito para um conjunto finito de estímulos, a equivalência formal compara matematicamente a implementação sintetizada com o design de referência. O Synopsys Formality realiza inicialmente o casamento dos pontos de comparação por meio do comando \texttt{match} e, em seguida, prova a equivalência lógica desses pontos com o comando \texttt{verify}.

O arquivo SVF (Setup/Verification File) foi utilizado como guidance entre a síntese e a verificação. Esse arquivo registra as transformações realizadas pelo Design Compiler e permite que o Formality interprete alterações como otimizações lógicas, mudanças de nomes e remoção de fronteiras hierárquicas. Foram executadas duas rodadas independentes: uma com \texttt{compile\_ultra -no\_autoungroup}, preservando a hierarquia, e outra com \texttt{compile\_ultra}, permitindo o desagrupamento automático.

% ----------------------------------------------------------
\section{Objetivos}
% ----------------------------------------------------------

Os objetivos específicos da atividade foram:

\begin{itemize}[leftmargin=1.2cm,itemsep=2pt]
    \item gerar um SVF próprio para cada rodada de síntese;
    \item carregar o RTL como design de referência e a netlist como implementação;
    \item habilitar o modo \texttt{synopsys\_auto\_setup};
    \item executar \texttt{match} e \texttt{verify};
    \item analisar operações SVF aceitas e rejeitadas;
    \item verificar a existência de pontos passing, failing e unmatched;
    \item comparar os fluxos grouped e ungrouped;
    \item documentar o sign-off formal com evidências reproduzíveis.
\end{itemize}

% ----------------------------------------------------------
\section{Fundamentos da Equivalência Formal}
% ----------------------------------------------------------

O \textit{golden design}, ou design de referência, corresponde ao RTL SystemVerilog original. Nesta atividade, ele é composto pelos módulos \texttt{vending\_pkg}, \texttt{credit\_reg}, \texttt{memory}, \texttt{comparator}, \texttt{subtractor}, \texttt{control\_unit} e \texttt{vending\_top}. O \textit{revised design}, ou implementação, corresponde à netlist Verilog mapeada gerada após \texttt{compile\_ultra}.

Os pontos de comparação incluem saídas primárias e elementos de estado. O comando \texttt{match} estabelece correspondências entre os pontos do RTL e da netlist, enquanto \texttt{verify} demonstra se a lógica entre cada par é equivalente. Um resultado \texttt{SUCCEEDED} indica que todos os pontos relevantes foram provados equivalentes. Pontos \texttt{failing} representam divergências lógicas, e pontos \texttt{unmatched} representam elementos que não encontraram correspondência entre os dois designs.

O SVF é central no fluxo porque mantém um trilho de auditoria das transformações da síntese. Cada netlist foi verificada com o SVF produzido na mesma execução de \texttt{compile\_ultra}, evitando a mistura de artefatos de rodadas distintas.

% ----------------------------------------------------------
\section{Ambiente e Organização dos Arquivos}
% ----------------------------------------------------------

As execuções foram realizadas no servidor \texttt{srv-microeletronica3}. As versões utilizadas foram:

\begin{table}[H]
\centering
\caption{Ferramentas utilizadas}
\label{tab:ferramentas}
\begin{tabular}{ll}
\toprule
\textbf{Ferramenta} & \textbf{Versão} \\
\midrule
Synopsys Design Compiler & X-2025.06-SP2 \\
Synopsys Formality & X-2025.06-SP3 \\
Biblioteca tecnológica & SAED32 RVT, \texttt{saed32rvt\_tt1p05v25c.db} \\
\bottomrule
\end{tabular}
\end{table}

A atividade foi organizada em diretórios separados para RTL, síntese, Formality, logs e relatórios. Os artefatos grouped e ungrouped receberam nomes independentes, incluindo netlists, SVFs e pastas de reports.

\begin{lstlisting}[style=reportstyle,caption={Estrutura principal da Atividade 02.}]
atividade_02_formality/
|-- rtl/
|-- libs/
|-- synth/
|   |-- synth_grouped.tcl
|   |-- synth_ungrouped.tcl
|   |-- netlist/
|   |-- reports/
|   `-- logs/
`-- fm/
    |-- formality_grouped.tcl
    |-- formality_ungrouped.tcl
    |-- reports/grouped/
    |-- reports/ungrouped/
    `-- logs/
\end{lstlisting}

% ----------------------------------------------------------
\section{Metodologia}
% ----------------------------------------------------------

\subsection{Síntese grouped}

Na primeira rodada, \texttt{set\_svf} foi ativado antes da compilação e a opção \texttt{-no\_autoungroup} foi utilizada para preservar a hierarquia dos módulos. Ao final, o SVF foi encerrado com \texttt{set\_svf -off}, e uma netlist específica foi escrita.

\lstinputlisting[
    style=tclstyle,
    firstline=26,
    lastline=63,
    caption={Trecho principal de \texttt{synth\_grouped.tcl}.},
    label={lst:synth-grouped}
]{dados/synth_grouped.tcl}

\subsection{Síntese ungrouped}

Na segunda rodada, o comando \texttt{compile\_ultra} foi executado sem \texttt{-no\_autoungroup}. Isso permitiu que o Design Compiler removesse fronteiras hierárquicas quando considerou vantajoso. Um novo SVF e uma nova netlist foram produzidos para essa execução.

\lstinputlisting[
    style=tclstyle,
    firstline=24,
    lastline=62,
    caption={Trecho principal de \texttt{synth\_ungrouped.tcl}.},
    label={lst:synth-ungrouped}
]{dados/synth_ungrouped.tcl}

\imagemrelatorio
{01_scripts_sintese_grouped_ungrouped.png}
{0.98}
{Trechos dos scripts de síntese grouped e ungrouped, mostrando \texttt{set\_svf}, \texttt{compile\_ultra} e a escrita das netlists.}
{fig:scripts-sintese}

\subsection{Geração dos artefatos}

A Figura~\ref{fig:artefatos-gerados} comprova a criação de dois SVFs e duas netlists. O SVF grouped possui 17.152 bytes e a netlist grouped 29.480 bytes. O SVF ungrouped possui 18.944 bytes e a netlist ungrouped 27.757 bytes. Os tamanhos diferentes confirmam que os arquivos pertencem a rodadas distintas.

\imagemrelatorio
{02_svfs_e_netlists_gerados.png}
{0.98}
{SVFs e netlists gerados nas duas rodadas de síntese.}
{fig:artefatos-gerados}

\subsection{Configuração do Formality}

Em cada script, o modo \texttt{synopsys\_auto\_setup} foi habilitado, a biblioteca tecnológica foi carregada, o SVF correspondente foi associado e os designs foram lidos em containers separados. O RTL foi carregado como referência em \texttt{r:/WORK/vending\_top}, enquanto a netlist foi carregada como implementação em \texttt{i:/WORK/vending\_top}.

Na versão instalada, foi adotada a ordem funcional \texttt{set synopsys\_auto\_setup true} antes de \texttt{set\_svf}. O próprio relatório final confirmou que o modo automático estava habilitado durante as duas verificações.

\lstinputlisting[
    style=tclstyle,
    firstline=7,
    lastline=74,
    caption={Fluxo de leitura, match, verify e geração de reports na configuração grouped.},
    label={lst:fm-grouped}
]{dados/formality_grouped.tcl}

\imagemrelatorio
{09_configuracao_scripts_formality.png}
{0.98}
{Configuração dos scripts Formality para as rodadas grouped e ungrouped.}
{fig:config-formality}

Os scripts entregues são scripts manuais e incluem diretamente todos os comandos de setup, \texttt{match}, \texttt{verify} e geração de relatórios. Não foi necessário utilizar \texttt{match -certain} nem \texttt{set\_user\_match}, pois o casamento padrão guiado pelo SVF resolveu todos os pontos.

% ----------------------------------------------------------
\section{Resultados da Verificação}
% ----------------------------------------------------------

\subsection{Resultado grouped}

A execução grouped terminou com código de saída zero e veredito \texttt{Verification SUCCEEDED}. Foram provados 101 pontos de comparação equivalentes, sendo 21 portas e 80 flip-flops. Nenhum ponto foi classificado como failing.

\imagemrelatorio
{04_status_final_grouped.png}
{0.90}
{Status final da verificação formal para a netlist grouped.}
{fig:status-grouped}

\begin{table}[H]
\centering
\caption{Resumo da verificação grouped}
\label{tab:grouped}
\begin{tabular}{lr}
\toprule
\textbf{Métrica} & \textbf{Resultado} \\
\midrule
Veredito & \texttt{Verification SUCCEEDED} \\
Compare points passing & 101 \\
Portas & 21 \\
Flip-flops (DFF) & 80 \\
Failing points & 0 \\
Unmatched points & 0 \\
Operações SVF aceitas & 64 \\
Operações SVF rejeitadas & 0 \\
Tempo do Formality & 4 s \\
Memória máxima & 1032 MB \\
\bottomrule
\end{tabular}
\end{table}

\subsection{Resultado ungrouped}

A execução ungrouped também terminou com código de saída zero e veredito \texttt{Verification SUCCEEDED}. O total de pontos foi idêntico ao da execução grouped: 101 compare points, compostos por 21 portas e 80 flip-flops, sem failing points.

\imagemrelatorio
{05_status_final_ungrouped.png}
{0.90}
{Status final da verificação formal para a netlist ungrouped.}
{fig:status-ungrouped}

\begin{table}[H]
\centering
\caption{Resumo da verificação ungrouped}
\label{tab:ungrouped}
\begin{tabular}{lr}
\toprule
\textbf{Métrica} & \textbf{Resultado} \\
\midrule
Veredito & \texttt{Verification SUCCEEDED} \\
Compare points passing & 101 \\
Portas & 21 \\
Flip-flops (DFF) & 80 \\
Failing points & 0 \\
Unmatched points & 0 \\
Operações SVF aceitas & 69 \\
Operações SVF rejeitadas & 0 \\
Tempo do Formality & 4 s \\
Memória máxima & 1032 MB \\
\bottomrule
\end{tabular}
\end{table}

\subsection{Operações SVF aceitas e rejeitadas}

A configuração grouped apresentou 64 operações SVF aceitas e nenhuma rejeitada. A configuração ungrouped apresentou 69 operações aceitas e nenhuma rejeitada. Portanto, o guidance foi consumido sem rejeições nas duas rodadas.

\imagemrelatorio
{06_operacoes_svf_aceitas_rejeitadas.png}
{0.98}
{Contagem de operações SVF aceitas e rejeitadas nas duas verificações.}
{fig:svf-contagem}

A diferença de cinco operações aceitas é explicada pelas operações \texttt{guide\_ungroup} registradas para cinco instâncias do projeto:

\lstinputlisting[
    style=reportstyle,
    firstline=204,
    lastline=231,
    caption={Operações \texttt{guide\_ungroup} aceitas no SVF da síntese ungrouped.},
    label={lst:guide-ungroup}
]{dados/svf_accepted_ungrouped.rpt}

As instâncias desagrupadas foram \texttt{u\_control\_unit}, \texttt{u\_credit\_reg}, \texttt{u\_memory}, \texttt{u\_comparator} e \texttt{u\_subtractor}. O SVF permitiu que essas mudanças de hierarquia fossem interpretadas automaticamente pelo Formality.

Os arquivos \texttt{svf\_rejected.rpt} contêm apenas o cabeçalho do relatório, sem nenhuma entrada de operação rejeitada, o que confirma a contagem igual a zero.

\subsection{Failing e unmatched points}

Os relatórios finais das duas configurações registraram explicitamente \texttt{No failing compare points} e \texttt{No unmatched points}. Assim, não houve necessidade de executar \texttt{analyze\_points}, aplicar casamento manual ou justificar pontos não comparados.

\imagemrelatorio
{08a_grouped_sem_failing_unmatched.png}
{0.90}
{Relatórios grouped sem failing points e sem unmatched points.}
{fig:sem-problemas-grouped}

\imagemrelatorio
{08b_ungrouped_sem_failing_unmatched.png}
{0.90}
{Relatórios ungrouped sem failing points e sem unmatched points.}
{fig:sem-problemas-ungrouped}

% ----------------------------------------------------------
\section{Comparação entre Grouped e Ungrouped}
% ----------------------------------------------------------

A Tabela~\ref{tab:comparacao} resume as duas execuções. Ambas chegaram ao mesmo resultado funcional, mas a rodada ungrouped exigiu cinco operações SVF adicionais para representar o desagrupamento dos submódulos.

\begin{table}[H]
\centering
\caption{Comparação das verificações grouped e ungrouped}
\label{tab:comparacao}
\small
\begin{tabularx}{\textwidth}{>{\raggedright\arraybackslash}Xcc}
\toprule
\textbf{Métrica} & \textbf{Grouped} & \textbf{Ungrouped} \\
\midrule
Comando de síntese & \texttt{compile\_ultra -no\_autoungroup} & \texttt{compile\_ultra} \\
Tamanho da netlist & 29.480 bytes & 27.757 bytes \\
Tamanho do SVF & 17.152 bytes & 18.944 bytes \\
Veredito & SUCCEEDED & SUCCEEDED \\
Compare points passing & 101 & 101 \\
Portas & 21 & 21 \\
DFFs & 80 & 80 \\
Failing points & 0 & 0 \\
Unmatched points & 0 & 0 \\
Operações SVF aceitas & 64 & 69 \\
Operações SVF rejeitadas & 0 & 0 \\
Tempo do Formality & 4 s & 4 s \\
Memória máxima & 1032 MB & 1032 MB \\
Matching manual necessário & Não & Não \\
\bottomrule
\end{tabularx}
\end{table}

\imagemrelatorio
{10a_comparacao_final_grouped.png}
{0.98}
{Resumo final da execução grouped.}
{fig:comparacao-grouped}

\imagemrelatorio
{10b_comparacao_final_ungrouped.png}
{0.98}
{Resumo final da execução ungrouped.}
{fig:comparacao-ungrouped}

A netlist ungrouped ficou ligeiramente menor, enquanto o seu SVF ficou maior. Esse comportamento é coerente com o achatamento hierárquico: a implementação pode eliminar estruturas de módulo, mas o SVF precisa registrar transformações adicionais para permitir que o Formality reconstrua a correspondência lógica com o RTL hierárquico.

A igualdade do número de compare points demonstra que a remoção das fronteiras de módulo não alterou a quantidade de saídas e elementos de estado relevantes para a prova. O guidance eliminou a necessidade de \texttt{match -certain} ou de associações manuais, mesmo na configuração com hierarquia modificada.

% ----------------------------------------------------------
\section{Avisos e Análise de Consistência}
% ----------------------------------------------------------

Durante a síntese, o Design Compiler registrou avisos sobre um ramo \texttt{default} inalcançável no módulo \texttt{credit\_reg} e sobre a desconsideração de comandos \texttt{initial} no módulo \texttt{memory}. O Formality também informou que comandos \texttt{initial} não são suportados e apresentou avisos relacionados a células de potência da biblioteca SAED32.

Na rodada grouped, o log registrou ainda a existência de cinco nets não dirigidas na implementação. O modo \texttt{synopsys\_auto\_setup} configurou o tratamento de sinais não dirigidos de forma compatível com a síntese. Esses avisos não produziram failing points, unmatched points ou resultado inconclusivo. Como ambas as execuções terminaram com \texttt{Verification SUCCEEDED}, os avisos não invalidaram a equivalência provada para os compare points do design.

Também foi verificado que cada script possui exatamente dois comandos \texttt{set\_top}: um para o design de referência e outro para a implementação. Essa organização evita o erro de redefinição de top observado durante a fase inicial de depuração.

% ----------------------------------------------------------
\section{Discussão}
% ----------------------------------------------------------

O resultado demonstra que a netlist produzida pelo Design Compiler é sequencialmente equivalente ao RTL da vending machine nas duas estratégias de síntese. A preservação da hierarquia não foi necessária para que o Formality alcançasse o sign-off, pois o SVF foi capaz de orientar o casamento também na versão ungrouped.

O principal efeito da remoção de \texttt{-no\_autoungroup} apareceu no conteúdo do SVF: foram registradas e aceitas cinco operações \texttt{guide\_ungroup}, exatamente uma para cada submódulo achatado. Apesar dessa diferença estrutural, os resultados finais permaneceram idênticos: 101 pontos equivalentes, zero failing e zero unmatched.

Sem guidance, alterações de hierarquia e nomes poderiam exigir uma segunda passada de \texttt{match -certain} ou comandos \texttt{set\_user\_match}. Neste fluxo, nenhuma dessas intervenções foi necessária. Portanto, o SVF reduziu o trabalho manual de matching e forneceu um trilho contínuo entre a síntese e a prova formal.

% ----------------------------------------------------------
\section{Conclusão}
% ----------------------------------------------------------

A atividade atingiu o objetivo de fechar o sign-off formal da síntese do controlador de vending machine. Foram geradas duas netlists e dois arquivos SVF independentes, correspondentes às configurações grouped e ungrouped. Os scripts do Formality carregaram corretamente a biblioteca, o guidance, o RTL e as netlists, executando \texttt{match} e \texttt{verify} sem erros fatais.

As duas verificações terminaram com \texttt{Verification SUCCEEDED}. Em cada rodada, 101 compare points foram provados equivalentes, sendo 21 portas e 80 flip-flops. Não houve failing points, unmatched points ou operações SVF rejeitadas. A configuração ungrouped apresentou cinco operações aceitas adicionais, associadas ao desagrupamento dos módulos, mas manteve o mesmo resultado funcional e o mesmo tempo de prova.

Conclui-se que as otimizações realizadas pelo Design Compiler, inclusive as mudanças de hierarquia da síntese ungrouped, preservaram integralmente o comportamento do RTL original. O uso do SVF foi suficiente para eliminar matching manual e documentar de forma reproduzível a equivalência entre o design de referência e as implementações sintetizadas.

% ----------------------------------------------------------
\section{Referências}
% ----------------------------------------------------------

\begin{itemize}[leftmargin=1.2cm,itemsep=2pt]
    \item CI Expert. \textit{Roteiro de Verificação de Equivalência Formal (Formality) - Vending Machine Controller}.
    \item Synopsys. \textit{Design Compiler User Guide}.
    \item Synopsys. \textit{Formality User Guide}.
    \item Synopsys. \textit{Formality Automated Setup File (SVF) Manual}.
\end{itemize}

\appendix

% ----------------------------------------------------------
\section{Resumo textual das execuções}
% ----------------------------------------------------------

O arquivo abaixo foi gerado diretamente no servidor a partir dos reports e logs das duas verificações.

\lstinputlisting[
    style=reportstyle,
    caption={Resumo consolidado das execuções do Formality.}
]{dados/evidencias_formality.txt}

% ----------------------------------------------------------
\section{Arquivos de evidência incluídos}
% ----------------------------------------------------------

O pacote LaTeX acompanha os scripts e relatórios utilizados na elaboração deste documento:

\begin{itemize}[leftmargin=1.2cm,itemsep=2pt]
    \item \texttt{dados/synth\_grouped.tcl} e \texttt{dados/synth\_ungrouped.tcl};
    \item \texttt{dados/formality\_grouped.tcl} e \texttt{dados/formality\_ungrouped.tcl};
    \item reports finais de status, failing e unmatched;
    \item reports de operações SVF aceitas e rejeitadas;
    \item imagens reais obtidas no servidor durante a execução.
\end{itemize}

\end{document}
